\documentclass[11pt,a4paper]{article}

\usepackage[french]{babel}
\usepackage[utf8]{inputenc}
\usepackage[T1]{fontenc}
\usepackage[margin=2.5cm]{geometry}
\usepackage{titlesec}
\usepackage{enumitem}
\usepackage{longtable}
\usepackage{array}
\usepackage{xcolor}
\usepackage{booktabs}
\usepackage{hyperref}
\usepackage{fancyhdr}
\usepackage{tabularx}

\definecolor{primaryblue}{RGB}{30,58,95}
\definecolor{accentgold}{RGB}{201,162,39}

\titleformat{\section}
  {\Large\bfseries\color{primaryblue}}{\thesection}{1em}{}
\titleformat{\subsection}
  {\large\bfseries\color{primaryblue}}{\thesubsection}{1em}{}
\titleformat{\subsubsection}
  {\normalsize\bfseries\color{accentgold}}{\thesubsubsection}{1em}{}

\pagestyle{fancy}
\fancyhf{}
\fancyhead[L]{Application de gestion d'évènements de mariage}
\fancyhead[R]{\thepage}
\renewcommand{\headrulewidth}{0.4pt}

\hypersetup{
  colorlinks=true,
  linkcolor=primaryblue,
  urlcolor=primaryblue,
  pdftitle={Cahier des charges - Application de gestion de mariages},
  pdfauthor={Michel Vaneck Kengne}
}

\newcolumntype{L}[1]{>{\raggedright\arraybackslash}p{#1}}

\title{\color{primaryblue}\textbf{Cahier des charges, conception et planification\\
Application de gestion d'évènements de mariage}}
\author{Michel Vaneck Kengne}
\date{\today}

\begin{document}

\maketitle
\thispagestyle{empty}

\vspace{1cm}
\noindent\textbf{Version :} 1.0 \\
\textbf{Statut :} Document de cadrage initial

\newpage
\tableofcontents
\newpage

% =====================================================
\section{Présentation du projet}
% =====================================================

\subsection{Contexte}
Le projet consiste à concevoir une application web permettant de créer, pour chaque
mariage, un site évènementiel personnalisé (invitations, programme, confirmation de
présence, gestion des tables, partage de photos et vidéos). Plutôt que de développer
un site unique par client, l'application repose sur un \textbf{moteur de gabarits
(template engine) configurable} : le code reste identique pour tous les mariages,
seule la configuration de chaque mariage change (contenu, sections actives, ordre,
couleurs, polices).

\subsection{Objectifs}
\begin{itemize}[leftmargin=1.2cm]
  \item Permettre à un administrateur de créer un compte pour chaque client (mariage)
        et de générer un site unique accessible via un lien.
  \item Permettre au client de personnaliser son site (contenu, apparence, sections
        actives, ordre d'affichage) sans toucher au code.
  \item Permettre aux invités de consulter le programme, confirmer leur présence,
        retrouver leur table et partager photos/vidéos.
  \item Fournir une expérience responsive et installable (PWA).
\end{itemize}

\subsection{Acteurs du système}
\begin{longtable}{L{3.5cm}L{10.5cm}}
\toprule
\textbf{Acteur} & \textbf{Rôle} \\
\midrule
\endhead
Administrateur & Crée et gère les comptes clients, active/désactive les liens, définit
la période d'accès. \\
Client (mariés) & Configure son site : contenu, apparence, sections, invités, tables,
galerie. \\
Invité & Consulte le site du mariage, confirme sa présence, retrouve sa table, partage
des médias, laisse un témoignage. \\
\bottomrule
\end{longtable}

% =====================================================
\section{Cahier des charges fonctionnel}
% =====================================================

\subsection{Module Administration (super-admin)}
\begin{itemize}[leftmargin=1.2cm]
  \item Création d'un compte client : nom, téléphone, type d'évènement, période
        d'accès, couleur, police, logo, choix du template (4 modèles de design :
        romantique floral, royal or, minimal épuré, héritage africain),
        génération automatique du lien du site.
  \item Activation / désactivation d'un compte client.
  \item Gestion de la période d'accès : définition des dates de disponibilité du lien,
        possibilité de désactiver le lien à tout moment.
\end{itemize}

\subsection{Module Compte client (mariés)}
\subsubsection{Configuration générale}
\begin{itemize}[leftmargin=1.2cm]
  \item Configuration de l'espace : ordre d'affichage des modules, activation ou
        désactivation de chaque module.
  \item Modules disponibles : photo de couverture (format mobile et desktop) ou
        vidéo, lieu (titre, adresse, image, localisation Google Maps), compte à
        rebours, image/vidéo, programme, partie mixte.
  \item Le programme peut comporter plusieurs étapes (dot, mairie, mariage religieux,
        église, soirée...), chacune avec sa propre date et son propre lieu.
  \item La « partie mixte » accepte des blocs de type texte, image, vidéo ou
        questionnaire, ajoutés et réordonnés librement (glisser-déposer).
\end{itemize}

\subsubsection{Gestion des invités}
\begin{itemize}[leftmargin=1.2cm]
  \item Ajout manuel d'un invité (nom, téléphone).
  \item Envoi du lien de l'évènement à l'invité, qui déroule ensuite le programme.
  \item Suivi de la présence par date et par sous-évènement (ex : présent à la
        mairie mais pas à la soirée), via case à cocher.
  \item Distinction des invités selon leurs réponses au questionnaire.
  \item Segmentation des invités par sous-évènement (dot / mairie / église / soirée).
\end{itemize}

\subsubsection{Gestion des tables}
\begin{itemize}[leftmargin=1.2cm]
  \item Création de tables : nom, capacité (optionnelle).
  \item Tables créées soit pour un sous-évènement précis (ex : tables de la soirée
        distinctes des tables de la mairie), soit communes à tout l'évènement.
  \item Affectation des invités aux tables, avec décompte automatique des places
        restantes à chaque ajout.
  \item Côté invité : recherche de sa table par son nom (« Find your table »).
\end{itemize}

\subsubsection{Galerie photos / vidéos}
\begin{itemize}[leftmargin=1.2cm]
  \item Upload d'images et de vidéos par l'administrateur du site (les mariés).
  \item Upload également possible par les invités ; les médias envoyés par les
        invités restent en attente jusqu'à validation par l'admin avant d'apparaître
        dans la galerie publique.
  \item Génération d'un QR code renvoyant vers la galerie, pour un partage rapide
        le jour de l'évènement.
\end{itemize}

\subsubsection{Témoignages}
\begin{itemize}[leftmargin=1.2cm]
  \item Formulaire permettant aux invités de laisser un témoignage.
  \item Possibilité de restreindre qui peut témoigner (accès limité aux invités
        identifiés / invités explicitement).
\end{itemize}

\subsection{Besoins non fonctionnels}
\begin{itemize}[leftmargin=1.2cm]
  \item \textbf{Responsive \& PWA} : interface adaptée mobile/desktop, installable
        sur l'écran d'accueil, chargement rapide, cache pour certaines
        fonctionnalités hors-ligne.
  \item \textbf{Multi-tenant} : un seul code applicatif sert un nombre illimité de
        mariages, chacun isolé par sa configuration et ses données.
  \item \textbf{Sécurité} : isolement des données par mariage (Row Level Security),
        accès admin protégé, liens invités non devinables.
  \item \textbf{Performance} : chargement rapide de la galerie et des médias
        (compression, chargement différé).
\end{itemize}

% =====================================================
\section{Conception}
% =====================================================

\subsection{Architecture générale}
L'application est structurée en trois couches~:

\begin{enumerate}[leftmargin=1.2cm]
  \item \textbf{Panel Super-Admin} : gestion des comptes clients (mariages).
  \item \textbf{Panel Client} : configuration du site (sections, apparence,
        invités, tables, galerie, témoignages).
  \item \textbf{Site public du mariage} : rendu du site accessible via
        \texttt{/w/\{slug\}}, généré à partir de la configuration du mariage.
\end{enumerate}

Le moteur de rendu (template engine) lit la configuration d'un mariage, filtre les
sections actives, les trie par position, puis affiche les composants correspondants.
Aucune section n'est codée en dur pour un mariage particulier : tout passe par la
configuration stockée en base.

\subsection{Choix techniques}
\begin{longtable}{L{4cm}L{10cm}}
\toprule
\textbf{Couche} & \textbf{Choix} \\
\midrule
\endhead
Frontend & React + TypeScript + Vite, Tailwind CSS, librairie de drag \& drop,
PWA / Service Worker \\
Backend / Données & Supabase (PostgreSQL, authentification, stockage de fichiers,
Row Level Security) \\
Partage & QR codes générés dynamiquement (galerie, RSVP, localisation) \\
\bottomrule
\end{longtable}

\subsection{Modèle de données (synthèse)}
\begin{longtable}{L{3.5cm}L{10.5cm}}
\toprule
\textbf{Table} & \textbf{Rôle} \\
\midrule
\endhead
\texttt{weddings} & Un enregistrement par mariage : slug, noms, template, thème,
période d'accès, statut actif. \\
\texttt{wedding\_events} & Sous-évènements d'un mariage (dot, mairie, église, soirée),
chacun avec sa date et son lieu. \\
\texttt{sections} & Modules affichables du site : type, activé/désactivé, position,
contenu (JSON). \\
\texttt{guests} & Invités d'un mariage, avec les sous-évènements auxquels ils sont
conviés. \\
\texttt{attendance} & Présence d'un invité pour un sous-évènement donné (présent /
absent / en attente), nombre d'accompagnants. \\
\texttt{guest\_responses} & Réponses des invités au questionnaire, utilisées pour
la segmentation. \\
\texttt{tables\_plan} & Tables, rattachées à un sous-évènement ou communes à tout
le mariage. \\
\texttt{table\_assignments} & Affectation des invités aux tables. \\
\texttt{media} & Photos / vidéos de la galerie, avec statut de modération. \\
\texttt{testimonials} & Témoignages des invités, avec statut de modération. \\
\bottomrule
\end{longtable}

% =====================================================
\section{Découpage du travail en phases}
% =====================================================

\subsection{Vue d'ensemble}
Le développement est découpé en huit phases, chacune livrant un incrément
fonctionnel testable.

\subsection{Phase 1 --- Socle technique}
\begin{itemize}[leftmargin=1.2cm]
  \item Initialisation du projet React + TypeScript + Vite + Tailwind.
  \item Mise en place du projet Supabase : base de données, authentification,
        stockage.
  \item Création du schéma de base de données (tables ci-dessus) et des règles
        Row Level Security.
  \item Mise en place du routage (\texttt{/admin}, \texttt{/app}, \texttt{/w/\{slug\}}).
\end{itemize}

\subsection{Phase 2 --- Panel Super-Admin}
\begin{itemize}[leftmargin=1.2cm]
  \item Authentification super-admin.
  \item Création / édition / désactivation d'un compte client (mariage).
  \item Génération automatique du lien et du slug.
  \item Gestion de la période d'accès.
\end{itemize}

\subsection{Phase 3 --- Moteur de gabarits (template engine)}
\begin{itemize}[leftmargin=1.2cm]
  \item Implémentation des composants de section (cover, lieu, countdown,
        programme, partie mixte, galerie, RSVP, tables, témoignages).
  \item Système de lecture de configuration : filtrage \texttt{enabled}, tri par
        \texttt{position}.
  \item Interface d'activation / désactivation et de réordonnancement
        (glisser-déposer) côté panel client.
  \item Système de thème (couleurs, polices, 4 modèles de design à part
        entière : palette, typographies et ornements décoratifs propres à
        chaque modèle).
\end{itemize}

\subsection{Phase 4 --- Programme et sous-évènements}
\begin{itemize}[leftmargin=1.2cm]
  \item Création des sous-évènements (dot, mairie, église, soirée) avec date et
        lieu.
  \item Affichage du programme multi-étapes côté site public.
  \item Intégration de la localisation (Google Maps).
\end{itemize}

\subsection{Phase 5 --- Invités et RSVP}
\begin{itemize}[leftmargin=1.2cm]
  \item Ajout manuel des invités (nom, téléphone) et association aux
        sous-évènements.
  \item Formulaire de confirmation de présence par sous-évènement.
  \item Questionnaire personnalisable et enregistrement des réponses.
  \item Tableau de bord de suivi (présents / absents / en attente) et segmentation
        des invités.
\end{itemize}

\subsection{Phase 6 --- Gestion des tables}
\begin{itemize}[leftmargin=1.2cm]
  \item Création des tables (nom, capacité), liées à un sous-évènement ou communes.
  \item Affectation des invités aux tables avec décompte des places restantes.
  \item Fonction « Find your table » côté invité.
\end{itemize}

\subsection{Phase 7 --- Galerie, témoignages et QR codes}
\begin{itemize}[leftmargin=1.2cm]
  \item Upload de médias par l'admin et par les invités.
  \item File de modération (validation avant publication).
  \item Génération des QR codes (galerie, RSVP, localisation).
  \item Formulaire de témoignages avec restriction d'accès optionnelle.
\end{itemize}

\subsection{Phase 8 --- PWA, tests et mise en production}
\begin{itemize}[leftmargin=1.2cm]
  \item Configuration PWA (manifest, service worker, installation, cache).
  \item Tests de bout en bout sur mobile et desktop.
  \item Optimisation des performances (images, chargement différé).
  \item Déploiement et mise à disposition du premier mariage pilote.
\end{itemize}

\subsection{Synthèse}
\begin{longtable}{L{1.5cm}L{6cm}L{6.5cm}}
\toprule
\textbf{Phase} & \textbf{Objectif} & \textbf{Livrable} \\
\midrule
\endhead
1 & Socle technique & Projet initialisé, base de données prête \\
2 & Panel Super-Admin & Création de comptes clients fonctionnelle \\
3 & Moteur de gabarits & Sections configurables et réordonnables \\
4 & Programme \& sous-évènements & Programme multi-étapes affiché \\
5 & Invités \& RSVP & Confirmation de présence et segmentation \\
6 & Gestion des tables & Affectation et recherche de table \\
7 & Galerie, témoignages, QR codes & Partage de médias et témoignages \\
8 & PWA \& mise en production & Application déployée, mariage pilote \\
\bottomrule
\end{longtable}

\end{document}
